% Generated draft: requires mathematical and source review.
\documentclass[11pt]{article}
\usepackage[margin=25mm]{geometry}
\usepackage{fontspec}
\setmainfont{Latin Modern Roman}
\usepackage{amsmath,amssymb,mathtools,unicode-math}
\setmathfont{Latin Modern Math}
\usepackage{enumitem,needspace,xurl,hyperref}
\setlength{\parindent}{0pt}
\setlength{\parskip}{5pt}
\setlength{\emergencystretch}{3em}
\newcommand{\N}{\mathbb{N}}
\newcommand{\Z}{\mathbb{Z}}
\newcommand{\Q}{\mathbb{Q}}
\newcommand{\R}{\mathbb{R}}
\newcommand{\C}{\mathbb{C}}
\newcommand{\F}{\mathbb{F}}
\newcommand{\A}{\mathbb{A}}
\newcommand{\PP}{\mathbb{P}}
\newcommand{\E}{\mathbb{E}}
\newcommand{\eps}{\varepsilon}
\newcommand{\dd}{\,\mathrm{d}}
\begin{document}
\section*{UnsolvedMath ID: 701466}
\textbf{Source problem code:} catalogue entry
\subsection{Definitions and mathematical statement}
Let \(G\) be a group, with no finiteness assumption on \(G\). Use the following standard group-theoretic conventions. The group \(G\) is \emph{locally graded} if every nontrivial finitely generated subgroup \(H\leq G\) has a proper subgroup \(U<H\) with \([H:U]<\infty\). It has \emph{finite rank} if there exists an integer \(r\geq 0\) such that every finitely generated subgroup of \(G\) can be generated by at most \(r\) elements.

For a group \(H\), define its derived series by \(H^{(0)}=H\) and \(H^{(n+1)}=[H^{(n)},H^{(n)}]\), where \([K,K]\) is the subgroup generated by all commutators \(x^{-1}y^{-1}xy\) with \(x,y\in K\). A group is \emph{soluble} if its derived series reaches the trivial subgroup after finitely many steps. A group \(L\) is \emph{locally soluble} if every finitely generated subgroup of \(L\) is soluble. Interpret \emph{almost locally soluble} as having a locally soluble subgroup of finite index.

\textbf{Problem (M. Dixon; described in the supplied source as a well-known problem).} Is every locally graded group of finite rank almost locally soluble? Explicitly, for every group \(G\) satisfying the first two definitions, must there exist \(L\leq G\) such that
\[
 [G:L]<\infty
 \quad\text{and}\quad
 \forall\,H\leq L\text{ finitely generated},\quad
 \exists\,d\in\Z_{\geq 0}\text{ such that }H^{(d)}=\{1\}?
\]
The integer \(d\) may depend on \(H\); no uniform bound on derived lengths is required. The supplied record does not define the technical terms, so the conventions stated above require confirmation against the original source. Its label \emph{open} is supplied registry metadata, not an independently verified mathematical assessment.
\subsection{Short English statement}
Does every locally graded group whose finitely generated subgroups have a uniform bound on their numbers of generators contain a locally soluble subgroup of finite index? Here locally soluble means that each finitely generated subgroup is soluble, with its own possible derived length. The technical conventions and supplied open status require human review.
\subsection{Sources}
\begin{itemize}
\item Local UnsolvedMath registry, supplied complete source record: permanent ID \(701466\), exact source problem code \texttt{catalogue entry}, title \emph{Kourovka Notebook Problem 18.37}. The record supplies the status \emph{open} and was updated on 2026-10-16.
\item \emph{The Kourovka Notebook (Unsolved Problems in Group Theory)}, Problem 18.37, 18th issue (2014); also identified in the supplied background with the 21st issue edition (2026), arXiv identifier \texttt{1401.0300v48}. The supplied collection description names E. I. Khukhro and V. D. Mazurov as editors. This bibliographic information is supplied attribution and has not been independently checked.
\item Proposed by M. Dixon, with proposed year 2014 in the supplied registry. The statement explicitly calls this a ``well-known problem''; this attribution does not establish original authorship of the question.
\end{itemize}
\end{document}
